% GENERATED FILE. Edit canonical lessons, then run npm run generate:books.
% canonical-source-hash: fnv1a64:5eb55cc9fa1ad8cd
% canonical-lessons: SA-C236-sisuh, SA-C236-kumarah, SA-C236-kumari, SA-C236-yuvatih, SA-C236-purusah

\chapter{Baby, Boy, Maiden, Young woman, Person}
\label{ch:sa-baby-boy-maiden-young-woman-person}
\hlchaptermodality{236}

\begin{chapteropening}
\textbf{By the end of this chapter:} \emph{I can say a baby, a boy, a maiden, a young woman and a person.}

Complete the last lesson of chapter 236: I can say a baby, a boy, a maiden, a young woman and a person.
\end{chapteropening}

\section[\texorpdfstring{\sk{शिशुः} (śiśuḥ)}{शिशुः (śiśuḥ)}]{\sk{शिशुः} (śiśuḥ) — a baby, an infant}

\label{lesson:SA-C236-sisuh}



\begin{quote}
Before the new one: say the Sanskrit for embraces, then the Sanskrit for kisses.
\end{quote}

\subsection*{You'll want to know: \sk{शिशुः}}
\textbf{\sk{शिशुः}} — \emph{śiśuḥ} — \textquotedblleft{}a baby, an infant\textquotedblright{}.

\begin{grammarlens}[title={What you've built}]
One more word for this chapter.
\end{grammarlens}

\subsection*{Your turn}
\begin{itemize}
\raggedright
  \item \emph{Say it:} \emph{śiśuḥ}
  \item \emph{Say it:} \emph{śiśuḥ}, once more
  \item \emph{Say it:} say \emph{śiśuḥ}
  \item \emph{Recall:} say the Sanskrit for embraces, then the Sanskrit for kisses, then say \emph{śiśuḥ} again
\end{itemize}

\begin{tcolorbox}[breakable,colback=teal!4,colframe=teal!35!black,title={Before you move on}]
What does \sk{शिशुः} mean? (\textquotedblleft{}A baby, an infant\textquotedblright{}.) Say it once more.
\end{tcolorbox}
\section[\texorpdfstring{\sk{कुमारः} (kumāraḥ)}{कुमारः (kumāraḥ)}]{\sk{कुमारः} (kumāraḥ) — a boy, a youth}

\label{lesson:SA-C236-kumarah}



\begin{quote}
Before the new one: say the Sanskrit for kisses, then the Sanskrit for a baby.
\end{quote}

\subsection*{You'll want to know: \sk{कुमारः}}
\textbf{\sk{कुमारः}} — \emph{kumāraḥ} — \textquotedblleft{}a boy, a youth\textquotedblright{}.

\begin{grammarlens}[title={What you've built}]
One more word for this chapter.
\end{grammarlens}

\subsection*{Your turn}
\begin{itemize}
\raggedright
  \item \emph{Say it:} \emph{kumāraḥ}
  \item \emph{Say it:} \emph{kumāraḥ}, once more
  \item \emph{Say it:} say \emph{kumāraḥ}
  \item \emph{Recall:} say the Sanskrit for kisses, then the Sanskrit for a baby, then say \emph{kumāraḥ} again
\end{itemize}

\begin{tcolorbox}[breakable,colback=teal!4,colframe=teal!35!black,title={Before you move on}]
What does \sk{कुमारः} mean? (\textquotedblleft{}A boy, a youth\textquotedblright{}.) Say it once more.
\end{tcolorbox}
\section[\texorpdfstring{\sk{कुमारी} (kumārī)}{कुमारी (kumārī)}]{\sk{कुमारी} (kumārī) — a maiden, a girl}

\label{lesson:SA-C236-kumari}



\begin{quote}
Before the new one: say the Sanskrit for a baby, then the Sanskrit for a boy.
\end{quote}

\subsection*{You'll want to know: \sk{कुमारी}}
\textbf{\sk{कुमारी}} — \emph{kumārī} — \textquotedblleft{}a maiden, a girl\textquotedblright{}.

\begin{grammarlens}[title={What you've built}]
One more word for this chapter.
\end{grammarlens}

\subsection*{Your turn}
\begin{itemize}
\raggedright
  \item \emph{Say it:} \emph{kumārī}
  \item \emph{Say it:} \emph{kumārī}, once more
  \item \emph{Say it:} say \emph{kumārī}
  \item \emph{Recall:} say the Sanskrit for a baby, then the Sanskrit for a boy, then say \emph{kumārī} again
\end{itemize}

\begin{tcolorbox}[breakable,colback=teal!4,colframe=teal!35!black,title={Before you move on}]
What does \sk{कुमारी} mean? (\textquotedblleft{}A maiden, a girl\textquotedblright{}.) Say it once more.
\end{tcolorbox}
\section[\texorpdfstring{\sk{युवतिः} (yuvatiḥ)}{युवतिः (yuvatiḥ)}]{\sk{युवतिः} (yuvatiḥ) — a young woman}

\label{lesson:SA-C236-yuvatih}



\begin{quote}
Before the new one: say the Sanskrit for a boy, then the Sanskrit for a maiden.
\end{quote}

\subsection*{You'll want to know: \sk{युवतिः}}
\textbf{\sk{युवतिः}} — \emph{yuvatiḥ} — \textquotedblleft{}a young woman\textquotedblright{}.

\begin{grammarlens}[title={What you've built}]
One more word for this chapter.
\end{grammarlens}

\subsection*{Your turn}
\begin{itemize}
\raggedright
  \item \emph{Say it:} \emph{yuvatiḥ}
  \item \emph{Say it:} \emph{yuvatiḥ}, once more
  \item \emph{Say it:} say \emph{yuvatiḥ}
  \item \emph{Recall:} say the Sanskrit for a boy, then the Sanskrit for a maiden, then say \emph{yuvatiḥ} again
\end{itemize}

\begin{tcolorbox}[breakable,colback=teal!4,colframe=teal!35!black,title={Before you move on}]
What does \sk{युवतिः} mean? (\textquotedblleft{}A young woman\textquotedblright{}.) Say it once more.
\end{tcolorbox}
\section[\texorpdfstring{\sk{पुरुषः} (puruṣaḥ)}{पुरुषः (puruṣaḥ)}]{\sk{पुरुषः} (puruṣaḥ) — a person, a man}

\label{lesson:SA-C236-purusah}



\begin{quote}
Before the new one: say the Sanskrit for a maiden, then the Sanskrit for a young woman.
\end{quote}

\subsection*{You'll want to know: \sk{पुरुषः}}
\textbf{\sk{पुरुषः}} — \emph{puruṣaḥ} — \textquotedblleft{}a person, a man\textquotedblright{}.

\begin{grammarlens}[title={What you've built}]
One more word for this chapter.
\end{grammarlens}

\subsection*{Your turn}
\begin{itemize}
\raggedright
  \item \emph{Say it:} \emph{puruṣaḥ}
  \item \emph{Say it:} \emph{puruṣaḥ}, once more
  \item \emph{Say it:} say \emph{puruṣaḥ}
  \item \emph{Recall:} say the Sanskrit for a maiden, then the Sanskrit for a young woman, then say \emph{puruṣaḥ} again
\end{itemize}

\begin{tcolorbox}[breakable,colback=teal!4,colframe=teal!35!black,title={Before you move on}]
What does \sk{पुरुषः} mean? (\textquotedblleft{}A person, a man\textquotedblright{}.) Say it once more.
\end{tcolorbox}
